\documentclass[a4paper,10pt]{report}

\def\packagepath{../../../../../preambule}   % path du package principal
\usepackage{\packagepath/preambule}    % utilisation du fichier de configuration

\def\level{BTS }              % Classe
\def\course{CIEL 1}           % Matière
\def\eval{Intégration}          % Titre de l'évaluation

\def\visibleornot{visible}    % visible or invisible
\def\documentpath{./Sources_Latex}
\def\date{C107 - TP1 }
\renewcommand{\arraystretch}{1}  % Ecart dans les tableaux

\def\ptsexoA{0}
\def\ptsexoB{0}

\def\ptstotal{\ptsexoA+\ptsexoB}

\usepackage{hyperref}

\usetikzlibrary{shapes, arrows, positioning}
\usetikzlibrary{shapes.geometric, arrows, positioning, decorations.pathreplacing}

\tikzstyle{etat} = [draw, rounded corners=15pt, minimum width=2.5cm, minimum height=1.2cm, text centered, font=\bfseries]
\tikzstyle{fleche} = [thick, ->, >=stealth]

\usepackage{float}
\usepackage[T1]{fontenc}

% --- L'INTERRUPTEUR ---
\booltrue{afficheCorrige} % Mettez \boolfalse pour masquer le texte (le cadre reste)
%\boolfalse{afficheCorrige} % Mettez \boolfalse pour masquer le texte (le cadre reste)

% --- LA COMMANDE DE CADRE FIXE ---
% #1: Largeur (ex: \linewidth)
% #2: Hauteur (ex: 3cm)
% #3: Contenu de la correction

%\cadreCorrection{width \linewidth}{cm}{}

\begin{document}





\renewcommand{\labelitemi}{\textbullet}

\pagestyle{DS_FP}
\NomPrenomNote{}


\bigskip



%%=============================================================
\textbf{PARTIE A: Mise en place de la fonction à étudier}


\medskip

L'objectif de cette partie est l'étude de la fonction $f$ définie sur $\R$ par:

\[f(x)=-0,5x^2+2x+1\]


\begin{enumerate}
    \item À l'aide de Géogebra, tracer la courbe représentative de la fonction $f$.\hfill $\square$
    \item Créer deux curseurs $a$ et $b$ compris entre 0 et 5 avec un pas de 0,1.\hfill $\square$
    \item Donner les valeurs suivantes à $a$ et $b$: $a=1$ et $b=4$.\hfill $\square$
    \end{enumerate}


    \bigskip

    \textbf{Partie B: Somme de Riemann (Approximation)}

    \medskip

    On souhaite diviser l'intervalle $[a\,;\,b]$ en $n$ sous-intervalles de même longueur. 

    \begin{enumerate}
        \item Créer un curseur $n$ compris entre 1 et 100 avec un pas de 1.\hfill $\square$
        \item Saisissez la commande \verb|S_inf = SommeInférieure(f, a, b, n)|\hfill $\square$
        \item Saisissez la commande \verb|S_sup = SommeSupérieure(f, a, b, n)|\hfill $\square$
        \item Changer la courleur de $S_{inf}$ et $S_{sup}$ pour les différencier.\hfill $\square$
        \item Observer les valeurs de $S_{inf}$ et $S_{sup}$ pour différentes valeurs de $n$. \hfill $\square$
        \item Que peut-on conjecturer sur la valeur de l'aire sous la courbe de $f$ entre $a$ et $b$ lorsque $n$ tend vers l'infini?\hfill $\square$
        
        \cadreCorrection{1}{3}{
                Lorsque $n$ tend vers l'infini, les sommes de Riemann $S_{inf}$ et $S_{sup}$ convergent vers la même valeur, qui correspond à l'aire sous la courbe de la fonction $f$ entre les points $a$ et $b$. On peut conjecturer que cette aire est égale à la limite de ces sommes lorsque $n$ tend vers l'infini.

                \[\lim_{n \to \infty} S_{\inf} = \lim_{n \to \infty} S_{\sup} = \lim_{n \to \infty} \int_{a}^{b} f(x) dx= \text{Aire sous la courbe de } f \text{ entre } a \text{ et } b\]        
                }       



    \end{enumerate}


    \bigskip

    \textbf{Partie C: Calcul de l'aire sous la courbe à l'aide de primitives}
    
    \medskip

    \begin{enumerate}
        \item Déterminer une primitive $F$ de la fonction $f$.\hfill $\square$
        
            \cadreCorrection{1}{3}{
                Pour trouver une primitive $F$ de la fonction $f(x) = -0,5x^2 + 2x + 1$, nous allons trouver une primitive de  chaque terme de la fonction séparément.

                \[F(x)=-0,5 \cdot \frac{x^3}{3} + 2 \cdot \frac{x^2}{2} + 1 \cdot x + C= -\frac{x^3}{6} + x^2 + x + C, \quad C \in \R\]
            }

        \item Calculer $F(b)-F(a)$ et interpréter ce résultat.\hfill $\square$
            \cadreCorrection{1}{3}{
                En calculant $F(b)-F(a)$, on obtient la valeur de l'aire sous la courbe de la fonction $f$ entre les points $a$ et $b$. 

                \[F(b)-F(a) = \int_{a}^{b} f(x) dx=\int_{1}^{4} f(x) dx= F(4)-F(1)=\frac{28}{3} - \frac{11}{6} = \frac{56 - 11}{6} = \frac{45}{6} = \frac{15}{2}=7,5\]
                Cette expression représente l'aire sous la courbe de $f$ entre les points $a$ et $b$ et au dessus de l'axe des abscisses.
            }

        \item Dans Géogebra, saisir la commande \verb|Aire = Intégrale(f, a, b)|. Comparer la valeur de $F(b)-F(a)$ avec celle de l'aire calculée dans Géogebra.\hfill $\square$
        
        
            \cadreCorrection{1}{3}{
                La valeur de $F(b)-F(a)$ est égale à l'aire calculée dans Géogebra, ce qui confirme que le calcul de l'aire sous la courbe de la fonction $f$ entre les points $a$ et $b$ est correct. 
            }

        \pagebreak
        \thispagestyle{DS_LP}
        \item Comparer la valeur de \verb|Aire| avec celles des sommes de Riemann $S_{inf}$ et $S_{sup}$ pour différentes valeurs de $n$.\hfill $\square$
        
            \cadreCorrection{1}{3}{
            
La valeur de \texttt{Aire} est égale à l'aire calculée dans Géogebra, qui correspond à l'aire sous la courbe de la fonction $f$ entre les points $a$ et $b$. En comparant cette valeur avec celles des sommes de Riemann $S_{inf}$ et $S_{sup}$ pour différentes valeurs de $n$, on peut observer que lorsque $n$ augmente, les valeurs de $S_{inf}$ et $S_{sup}$ se rapprochent de plus en plus de la valeur de \texttt{Aire}, confirmant ainsi la conjecture faite précédemment sur la convergence des sommes de Riemann vers l'aire sous la courbe.
             }


        \item Calculer la valeur moyenne de la fonction $f$ sur l'intervalle $[a\,;\,b]$ et interpréter ce résultat.\hfill $\square$
        
            \cadreCorrection{1}{3}{
                La valeur moyenne de la fonction $f$ sur l'intervalle $[a\,;\,b]$ est donnée par la formule :
                \[\mu_{[a; b]}=\frac{1}{b-a} \int_{a}^{b} f(x) dx = \frac{1}{4-1} \cdot \frac{15}{2} = \frac{1}{3} \cdot \frac{15}{2} = \frac{5}{2}=2,5\]
                Cette valeur correspond à la hauteur du rectangle dont l'aire est égale à l'aire sous la courbe de $f$ entre $a$ et $b$.
            }

        \item Créer avec Géogébra le rectangle $ABCD$ tel que $AB$ soit égal à la valeur moyenne de $f$ sur $[a\,;\,b]$ et que $AD$ soit égal à $b-a$. Vérifier que l'aire de ce rectangle est égale à l'aire sous la courbe de $f$ entre $a$ et $b$.\hfill $\square$
    
            \cadreCorrection{1}{4}{
                En créant le rectangle $ABCD$ avec $AB$ égal à la valeur moyenne de $f$ sur $[a\,;\,b]$ (soit 2,5) et $AD$ égal à $b-a$ (soit 3), on obtient un rectangle dont l'aire est donnée par :
                \[\text{Aire}_{ABCD} = AB \cdot AD = 2,5 \cdot 3 = 7,5\]
                Cette aire est égale à l'aire sous la courbe de $f$ entre les points $a$ et $b$, confirmant ainsi que le rectangle $ABCD$ a la même aire que celle sous la courbe de la fonction $f$ sur l'intervalle $[a\,;\,b]$.
            }

    \end{enumerate}



\end{document}